\chapter{The Rhythm Section: Lock Drums, Bass and Comping to the Same Beat}
\label{ch:rhythm}

\classquote{A rhythm section is not ``drums plus bass plus a guitar''. It is a three-part machine: one part sets the time, one pushes it forward, one adds colour.}{}{}

\section{What the rhythm section is responsible for}
\label{sec:rhythm-job}

The \term{rhythm section}{rhythm-section} is the group of voices that carries groove and root motion: usually drums, bass, rhythm guitar or keyboard, and percussion. It is responsible for at least four things:

\begin{enumerate}
\item \textbf{Set the time}: let the listener know how long a beat and a bar are. Once the time is stable, the melody can afford to be free.
\item \textbf{Set the groove}: decide where the accents fall, where space appears, and where the music pushes. This is the source of \term{groove}{groove}.
\item \textbf{Set the roots}: use bass and kick to establish the harmonic root and the low-frequency centre, so the harmony above has a foundation.
\item \textbf{Set the energy}: within one song, the density and velocity of the rhythm section decide whether a section walks or sprints.
\end{enumerate}

Listeners may forget a melody that is merely competent; they will always notice a rhythm section that is falling apart. So this chapter works in this order: make time and roots stable first, then talk about tricks.

\section{Three jobs: anchor, drive, colour}
\label{sec:rhythm-layers}

Dividing the rhythm section into three jobs is more useful than dividing it by instrument:

\begin{longtable}{p{0.16\textwidth}p{0.30\textwidth}p{0.46\textwidth}}
\toprule
Job & Usual members & Responsibility \\
\midrule
\endfirsthead
\toprule
Job & Usual members & Responsibility \\
\midrule
\endhead
Anchor & Kick, bass roots & Pins the weight and root of each bar; the listener uses it to decide where the beat is. \\
Drive & Snare, hi-hats, mid-range bass notes & Carries time from one beat to the next; decides whether the groove is even, pushing or clipped. \\
Colour & Percussion, rhythmic guitar/keyboard ornaments & Adds a different texture to the same groove; the section can stand without it, but style appears with it. \\
\bottomrule
\end{longtable}

The three jobs are not three fixed instrument groups. Bass can be only an anchor, or it can drive in the middle of the bar; a keyboard can be only colour, or it can carry harmony in the chorus. Before writing a beat, ask: who is setting time, who is pushing, who is adding colour? If all three want the same beat, the result is crowded.

\section{Drums: write the skeleton, then the density}
\label{sec:drums}

Drum writing starts with three things: kick, snare and hi-hat. Write their placement first; add fills and percussion later.

\begin{longtable}{p{0.20\textwidth}p{0.32\textwidth}p{0.40\textwidth}}
\toprule
Skeleton & Placement & Use and risk \\
\midrule
\endfirsthead
\toprule
Skeleton & Placement & Use and risk \\
\midrule
\endhead
Standard backbeat & Kick on beats 1 and 3; snare on 2 and 4 & The most stable pop/rock skeleton; risk: ordinary, so hats and fills must provide character. \\
Half-time & Kick on 1, snare on 3; the whole time feel widens & Ballads, slow rock, atmospheric tracks; risk: not enough drive, the chorus can drag. \\
Four-on-the-floor & Kick on every beat; snare on 2 and 4 or syncopated & Dance, electronic, disco; risk: the low end is always present and bass fights the kick. \\
Syncopated kick & Kick avoids the strong beats and lands on weak or off-beats & Funk, hip-hop, Latin; risk: the listener loses the beat unless an anchor remains stable. \\
\bottomrule
\end{longtable}

Follow two rules when writing drums:

\begin{itemize}
\item \textbf{Skeleton first, fills second.} The kick and snare placement must stand on its own within one bar; a fill is decoration, not structure.
\item \textbf{The snare is not the groove's lead.} The snare carries the backbeat and drive; the kick carries the weight. If they fight for the same beat, low and mid frequencies muddy together.
\end{itemize}

Hi-hats decide density: the verse can use quarters or eighths; the chorus can move to sixteenths. Do not keep one density from top to bottom. Percussion belongs to the colour job; the demo mutes it by default precisely to remind you that the skeleton must stand without it.

\section{Bass: roots first, groove second}
\label{sec:bass}

The bass's first job is to let the listener hear the harmonic root. Its second job is to make that root move. Reverse the order and the bass becomes a second melody.

\begin{longtable}{p{0.26\textwidth}p{0.64\textwidth}}
\toprule
Approach & Explanation \\
\midrule
\endfirsthead
\toprule
Approach & Explanation \\
\midrule
\endhead
Roots first & Return to the root every bar or half bar; the listener uses it to confirm the chord. \\
Same-note repetition & A practical line for stable groove: same-note repetition around 70\% or more; below 60\% tends to sound melodic and steal focus. \\
Limit leaps & Keep leaps of a fifth or more (5 semitones) under about 10\% of notes; passing phrases may break the rule, but do not spend the whole section jumping. \\
Register & A common starting range is E1--G2 (28--43). Frequent excursions above C3 empty the low centre and make the bass sound floaty. \\
Note length & Long notes are steadier; eighths push more. Use longer notes in the verse and shorter ones in the chorus; do not start with sixteenth-note bass. \\
\bottomrule
\end{longtable}

The bass-to-kick relationship matters more than the bass-to-snare relationship. There are two basic alignments:

\begin{itemize}
\item \textbf{Land together}: bass and kick hit the strong beat together; the low end is most stable and suits most pop/rock root beds.
\item \textbf{Offset}: the kick lands on the strong beat and the bass arrives half a beat later or on an off-beat; the groove pushes more and suits funk or electronic music, but needs a stronger anchor to avoid drifting.
\end{itemize}

\section{Comping: compress the chords, do not fill every beat}
\label{sec:comping}

Rhythmic guitar or keyboard \term{comping}{comping} is not ``hold the chord from the first beat to the last''. Its job is to give the chord at the right moments and leave space where it is not needed.

\begin{itemize}
\item \textbf{Give the harmonic skeleton}: decide which beats in each bar must make the chord change audible; the other beats can be empty or a single note.
\item \textbf{Avoid the bass root}: do not put the comping's lowest voice on the same root as the bass; keep the chord in the midrange and leave the low end to the bass.
\item \textbf{Leave a little space every four bars}: constant eighths or sixteenths tire the ear; one beat of space makes the drums and bass clearer.
\item \textbf{Choose strums or single notes}: full strums suit broad gestures; a single-note riff suits a heavier, steadier groove. Do not strum and fill on every beat.
\end{itemize}

\section{Aligning the three parts: three questions on one grid}
\label{sec:rhythm-align}

Put drums, bass and comping on the same eight-beat grid and check only three questions:

\begin{enumerate}
\item \textbf{Do kick and bass lock?} When they land together, is the low end too heavy? When they are offset, is there still an anchor?
\item \textbf{Do snare and comping cover each other?} When the snare lands on 2 and 4, the comping should not give a full chord on the same beat.
\item \textbf{Are hats and percussion too dense?} Density is relative to the section; sixteenths in the verse may be only the baseline in the chorus.
\end{enumerate}

\begin{figure}[htbp]
  \centering
  \begin{tikzpicture}[font=\small]
    \draw[step=1.2cm, gray!35, thin] (0,0) grid (9.6,3.6);
    \draw[thick] (0,0) rectangle (9.6,3.6);
    \node[anchor=east] at (-0.1,3.0) {kick};
    \node[anchor=east] at (-0.1,2.1) {snare};
    \node[anchor=east] at (-0.1,1.2) {bass};
    \node[anchor=east] at (-0.1,0.3) {comp};
    \foreach \x in {0,2.4,4.8,7.2} {
      \fill[yebanblue!75] (\x,2.75) rectangle ++(1.2,0.5);
    }
    \foreach \x in {1.2,3.6,6.0,8.4} {
      \fill[gray!70] (\x,1.85) rectangle ++(1.2,0.5);
    }
    \foreach \x in {0,1.2,2.4,3.6,4.8,6.0,7.2,8.4} {
      \fill[yebanblue!30] (\x,0.95) rectangle ++(1.1,0.45);
    }
    \foreach \x in {0,2.4,4.8,7.2} {
      \fill[gray!35] (\x,0.05) rectangle ++(1.2,0.45);
    }
    \node[font=\scriptsize, anchor=west] at (9.7,3.0) {strong beat};
    \node[font=\scriptsize, anchor=west] at (9.7,2.1) {backbeat};
    \node[font=\scriptsize, anchor=west] at (9.7,1.2) {root};
    \node[font=\scriptsize, anchor=west] at (9.7,0.3) {space};
  \end{tikzpicture}
  \caption{An eight-beat rhythm-section sketch: kick sets the strong beat, snare takes the backbeat, bass guards the root, comping gives chords on key beats and leaves space elsewhere.}
  \label{fig:rhythm-grid}
\end{figure}

The figure is not a pattern to copy; it is a way to point at ``who is responsible for this beat''. If all four rows land on every beat, the listener does not hear richness; they hear crowding.

\section{Realising it in Yeban}
\label{sec:rhythm-yeban}

\begin{enumerate}
\item \textbf{See the demo roles first.} In the demo project, drums are a MIDI track, bass is an audio track and percussion is an audio track muted by default. When listening to the finished groove, solo bass and drums to see how they lock. When writing your own bass, create a MIDI synth track so you can edit notes, velocity and note length.
\item \textbf{Create tracks.} Use ``Track $\rightarrow$ + New MIDI synth track (Ctrl+N)'' for drums, bass and comping; put audio material on audio tracks. Do not stuff drums, bass and comping into one track.
\item \textbf{Write the skeleton in the piano roll.} Write kick and snare placement first, then bass roots, then comping chords. Reverse the order and the ornaments will lead you.
\item \textbf{Quantise and humanise.} Use ``Edit $\rightarrow$ Quantize notes (1/16)'' to clean hand-played timing; then use ``Edit $\rightarrow$ Humanize notes'' for reproducible micro-differences in timing and velocity. Do not quantise the whole song into a straight line.
\item \textbf{Use the velocity lane.} Draw velocity note by note below the piano roll; do not use one velocity for the whole song. A common starting point is 90--105 in the verse and 105--120 in the chorus; it is not a law.
\item \textbf{Check with Solo/Mute.} Solo drums and bass to hear whether the low end locks; mute comping to hear whether the skeleton still stands; mute percussion to hear whether the colour layer is only decoration.
\item \textbf{Watch the low end in the mixer.} Press X / F8 to open the mixer and see whether kick and bass levels push against each other. Solve ``two parts fighting for the same beat and low end'' at the arranging stage; do not wait for EQ to cut it apart.
\end{enumerate}

\screenshot{arrangement-rhythm-section-groove}{The eight-beat skeleton across drum, bass and comping tracks, with mixer low-end levels: see who sets the strong beat, who takes the backbeat, and who leaves space}

\section{Common mistakes and fixes}
\label{sec:rhythm-mistakes}

\begin{longtable}{p{0.34\textwidth}p{0.58\textwidth}}
\toprule
Common mistake & Fix \\
\midrule
\endfirsthead
\toprule
Common mistake & Fix \\
\midrule
\endhead
The bass becomes a second melody & Return to roots; raise same-note repetition and keep leaps under 10\% so the listener hears the harmony first. \\
Kick and bass land on every beat together & Change to one strong/one weak hit or offset by half a beat; constant double low end becomes a rumble. \\
Snare and comping give full chords on the same beat & Leave space or use a single note at the snare; give the backbeat to the snare. \\
Hi-hats are sixteenths from top to bottom & Change density by section: eighths in the verse, sixteenths in the chorus; use space to create drive. \\
Everything is quantised equally tight & Align the skeleton to the grid, humanise ornaments and velocity; do not put every track on the same grid. \\
One flat velocity for the whole song & Leave dynamic levels between verse and chorus; the rhythm section also needs to breathe. \\
Comping's low voice fights the bass root & Remove or raise the comping's lowest voice; leave the low end to the bass. \\
Percussion runs from start to finish & Percussion is a colour layer; the skeleton must stand without it. Add and remove it by section. \\
\bottomrule
\end{longtable}

\section{Exercise}
\label{sec:rhythm-exercise}

\begin{exercise}
Build an eight-bar rhythm section and make three versions; record the differences:
\begin{enumerate}
\item \textbf{Standard}: kick on 1 and 3, snare on 2 and 4, bass mostly roots, comping gives a chord every two beats.
\item \textbf{Driving}: keep the skeleton, move the chorus hi-hats to sixteenths, and add one passing note on beat 3 of each bar in the bass. Record the same-note repetition rate and leap share.
\item \textbf{Half-time}: kick on 1, snare on 3, longer bass notes; compare the low-end centre and the ``walks/drags'' feel of all three versions.
\end{enumerate}
For each version, check with Solo/Mute: does the skeleton stand with drums and bass soloed? Does the groove survive when comping is muted? Is percussion only decoration when muted?
\end{exercise}

\section{Chapter takeaways}
\label{sec:rhythm-takeaways}

\begin{itemize}
\item The rhythm section sets time, groove, roots and energy.
\item Three jobs: anchor, drive, colour; for every beat, ask who is responsible.
\item Write the drum skeleton before the density; the snare is not the groove's lead.
\item Bass roots first; same-note repetition around 70\% or more and leaps under 10\% are practical lines for a stable groove.
\item Compress the chords, leave space and avoid the bass root; do not fill every beat.
\item Solve same-beat low-end conflicts at the arranging stage; do not leave them for EQ and compression.
\end{itemize}

\begin{helpmapping}
\begin{longtable}{p{0.45\textwidth}p{0.45\textwidth}}
\toprule
Section & Help anchor \\
\midrule
What the rhythm section is responsible for & \texttt{arrangement-rhythm-job} \\
Three jobs: anchor, drive, colour & \texttt{arrangement-rhythm-layers} \\
Drums: skeleton, then density & \texttt{arrangement-drums} \\
Bass: roots first & \texttt{arrangement-bass} \\
Comping: compress the chords & \texttt{arrangement-comping} \\
Aligning the three parts & \texttt{arrangement-rhythm-align} \\
Realising it in Yeban & \texttt{arrangement-rhythm-yeban} \\
Common mistakes and fixes & \texttt{arrangement-rhythm-mistakes} \\
Exercise & \texttt{arrangement-rhythm-exercise} \\
\bottomrule
\end{longtable}
\end{helpmapping}
